% ============================================================
% fig_som_proceso.tex: Proceso de entrenamiento SOM (pasos)
% Uso: % ============================================================
% fig_som_proceso.tex: Proceso de entrenamiento SOM (pasos)
% Uso: % ============================================================
% fig_som_proceso.tex: Proceso de entrenamiento SOM (pasos)
% Uso: % ============================================================
% fig_som_proceso.tex: Proceso de entrenamiento SOM (pasos)
% Uso: \input{figures/fig_som_proceso}
% ============================================================
\begin{figure}[H]
\centering
\begin{tikzpicture}[
    neuron/.style={circle, draw=blue!50, fill=blue!8, minimum size=0.55cm},
    bmu/.style={circle, draw=red!70, fill=red!20, minimum size=0.55cm, very thick},
    neighbor/.style={circle, draw=orange!70, fill=orange!10, minimum size=0.55cm, thick},
    input/.style={circle, draw=green!60!black, fill=green!15, minimum size=0.6cm, font=\tiny},
    arrow/.style={-{Stealth}, thick},
    label/.style={font=\small\bfseries, gray}
]

% ─── Grilla 3×3 en t=0 (aleatoria) ──────────────────────────
\node[label] at (1.0, 4.1) {$t = 0$ (inicial)};
\foreach \row in {0,1,2}{
    \foreach \col in {0,1,2}{
        \node[neuron] at (\col*0.9, 3.2-\row*0.9) {};
    }
}
\node[font=\tiny, gray] at (1.0, 0.8) {pesos aleatorios};

% ─── Flecha paso 1 ──────────────────────────────────────────
\draw[arrow, blue!60] (2.4, 2.2) -- (3.2, 2.2);
\node[font=\tiny, align=center, blue!70] at (2.8, 2.55)
    {1. presentar\\entrada $\bm{x}$};

% ─── Grilla con BMU marcada ──────────────────────────────────
\node[label] at (5.0, 4.1) {BMU encontrada};
% Entrada
\node[input] at (5.0, 4.55) {$\bm{x}$};
\foreach \row in {0,1,2}{
    \foreach \col in {0,1,2}{
        \pgfmathparse{int(\row==1 && \col==1)}
        \ifnum\pgfmathresult=1
            \node[bmu]      at (3.6+\col*0.9, 3.2-\row*0.9) {};
        \else
            \pgfmathparse{int(abs(\row-1)<=1 && abs(\col-1)<=1)}
            \ifnum\pgfmathresult=1
                \node[neighbor] at (3.6+\col*0.9, 3.2-\row*0.9) {};
            \else
                \node[neuron]   at (3.6+\col*0.9, 3.2-\row*0.9) {};
            \fi
        \fi
    }
}
\draw[red!60, dashed] (5.0, 4.4) -- (4.5, 2.3);
\node[font=\tiny, gray] at (5.0, 0.8) {BMU=rojo, vecinos=naranja};

% ─── Flecha paso 2 ──────────────────────────────────────────
\draw[arrow, blue!60] (6.4, 2.2) -- (7.2, 2.2);
\node[font=\tiny, align=center, blue!70] at (6.8, 2.55)
    {2. actualizar\\BMU+vecinos};

% ─── Grilla actualizada ─────────────────────────────────────
\node[label] at (9.0, 4.1) {pesos actualizados};
\foreach \row in {0,1,2}{
    \foreach \col in {0,1,2}{
        \pgfmathparse{int(\row==1 && \col==1)}
        \ifnum\pgfmathresult=1
            \node[bmu]      at (7.6+\col*0.9, 3.2-\row*0.9) {};
        \else
            \pgfmathparse{int(abs(\row-1)<=1 && abs(\col-1)<=1)}
            \ifnum\pgfmathresult=1
                \node[neighbor] at (7.6+\col*0.9, 3.2-\row*0.9) {};
            \else
                \node[neuron]   at (7.6+\col*0.9, 3.2-\row*0.9) {};
            \fi
        \fi
    }
}
\node[font=\tiny, gray, align=center] at (9.0, 0.8) {BMU y vecinos se\\acercan a $\bm{x}$};

% ─── Leyenda regla de actualización ─────────────────────────
\node[draw=blue!20, fill=blue!5, rounded corners, font=\footnotesize,
      inner sep=5pt, align=left] at (5.5, -0.5)
      {$\bm{w}_i(t{+}1) = \bm{w}_i(t) + \alpha(t)\, h_{ci}(t)\,[\bm{x}(t) - \bm{w}_i(t)]$};

\end{tikzpicture}
\caption{Proceso de entrenamiento de un SOM en una iteración. (1) Se presenta
una entrada $\bm{x}$ y se identifica la neurona ganadora o \textit{Best Matching
Unit} (BMU, en rojo), la de menor distancia a $\bm{x}$. (2) La BMU y sus
vecinas (naranja) actualizan sus pesos acercándose a $\bm{x}$ según la regla
mostrada, donde $\alpha(t)$ es la tasa de aprendizaje y $h_{ci}(t)$ es la
función de vecindad.}
\label{fig:som_proceso}
\end{figure}

% ============================================================
\begin{figure}[H]
\centering
\begin{tikzpicture}[
    neuron/.style={circle, draw=blue!50, fill=blue!8, minimum size=0.55cm},
    bmu/.style={circle, draw=red!70, fill=red!20, minimum size=0.55cm, very thick},
    neighbor/.style={circle, draw=orange!70, fill=orange!10, minimum size=0.55cm, thick},
    input/.style={circle, draw=green!60!black, fill=green!15, minimum size=0.6cm, font=\tiny},
    arrow/.style={-{Stealth}, thick},
    label/.style={font=\small\bfseries, gray}
]

% ─── Grilla 3×3 en t=0 (aleatoria) ──────────────────────────
\node[label] at (1.0, 4.1) {$t = 0$ (inicial)};
\foreach \row in {0,1,2}{
    \foreach \col in {0,1,2}{
        \node[neuron] at (\col*0.9, 3.2-\row*0.9) {};
    }
}
\node[font=\tiny, gray] at (1.0, 0.8) {pesos aleatorios};

% ─── Flecha paso 1 ──────────────────────────────────────────
\draw[arrow, blue!60] (2.4, 2.2) -- (3.2, 2.2);
\node[font=\tiny, align=center, blue!70] at (2.8, 2.55)
    {1. presentar\\entrada $\bm{x}$};

% ─── Grilla con BMU marcada ──────────────────────────────────
\node[label] at (5.0, 4.1) {BMU encontrada};
% Entrada
\node[input] at (5.0, 4.55) {$\bm{x}$};
\foreach \row in {0,1,2}{
    \foreach \col in {0,1,2}{
        \pgfmathparse{int(\row==1 && \col==1)}
        \ifnum\pgfmathresult=1
            \node[bmu]      at (3.6+\col*0.9, 3.2-\row*0.9) {};
        \else
            \pgfmathparse{int(abs(\row-1)<=1 && abs(\col-1)<=1)}
            \ifnum\pgfmathresult=1
                \node[neighbor] at (3.6+\col*0.9, 3.2-\row*0.9) {};
            \else
                \node[neuron]   at (3.6+\col*0.9, 3.2-\row*0.9) {};
            \fi
        \fi
    }
}
\draw[red!60, dashed] (5.0, 4.4) -- (4.5, 2.3);
\node[font=\tiny, gray] at (5.0, 0.8) {BMU=rojo, vecinos=naranja};

% ─── Flecha paso 2 ──────────────────────────────────────────
\draw[arrow, blue!60] (6.4, 2.2) -- (7.2, 2.2);
\node[font=\tiny, align=center, blue!70] at (6.8, 2.55)
    {2. actualizar\\BMU+vecinos};

% ─── Grilla actualizada ─────────────────────────────────────
\node[label] at (9.0, 4.1) {pesos actualizados};
\foreach \row in {0,1,2}{
    \foreach \col in {0,1,2}{
        \pgfmathparse{int(\row==1 && \col==1)}
        \ifnum\pgfmathresult=1
            \node[bmu]      at (7.6+\col*0.9, 3.2-\row*0.9) {};
        \else
            \pgfmathparse{int(abs(\row-1)<=1 && abs(\col-1)<=1)}
            \ifnum\pgfmathresult=1
                \node[neighbor] at (7.6+\col*0.9, 3.2-\row*0.9) {};
            \else
                \node[neuron]   at (7.6+\col*0.9, 3.2-\row*0.9) {};
            \fi
        \fi
    }
}
\node[font=\tiny, gray, align=center] at (9.0, 0.8) {BMU y vecinos se\\acercan a $\bm{x}$};

% ─── Leyenda regla de actualización ─────────────────────────
\node[draw=blue!20, fill=blue!5, rounded corners, font=\footnotesize,
      inner sep=5pt, align=left] at (5.5, -0.5)
      {$\bm{w}_i(t{+}1) = \bm{w}_i(t) + \alpha(t)\, h_{ci}(t)\,[\bm{x}(t) - \bm{w}_i(t)]$};

\end{tikzpicture}
\caption{Proceso de entrenamiento de un SOM en una iteración. (1) Se presenta
una entrada $\bm{x}$ y se identifica la neurona ganadora o \textit{Best Matching
Unit} (BMU, en rojo), la de menor distancia a $\bm{x}$. (2) La BMU y sus
vecinas (naranja) actualizan sus pesos acercándose a $\bm{x}$ según la regla
mostrada, donde $\alpha(t)$ es la tasa de aprendizaje y $h_{ci}(t)$ es la
función de vecindad.}
\label{fig:som_proceso}
\end{figure}

% ============================================================
\begin{figure}[H]
\centering
\begin{tikzpicture}[
    neuron/.style={circle, draw=blue!50, fill=blue!8, minimum size=0.55cm},
    bmu/.style={circle, draw=red!70, fill=red!20, minimum size=0.55cm, very thick},
    neighbor/.style={circle, draw=orange!70, fill=orange!10, minimum size=0.55cm, thick},
    input/.style={circle, draw=green!60!black, fill=green!15, minimum size=0.6cm, font=\tiny},
    arrow/.style={-{Stealth}, thick},
    label/.style={font=\small\bfseries, gray}
]

% ─── Grilla 3×3 en t=0 (aleatoria) ──────────────────────────
\node[label] at (1.0, 4.1) {$t = 0$ (inicial)};
\foreach \row in {0,1,2}{
    \foreach \col in {0,1,2}{
        \node[neuron] at (\col*0.9, 3.2-\row*0.9) {};
    }
}
\node[font=\tiny, gray] at (1.0, 0.8) {pesos aleatorios};

% ─── Flecha paso 1 ──────────────────────────────────────────
\draw[arrow, blue!60] (2.4, 2.2) -- (3.2, 2.2);
\node[font=\tiny, align=center, blue!70] at (2.8, 2.55)
    {1. presentar\\entrada $\bm{x}$};

% ─── Grilla con BMU marcada ──────────────────────────────────
\node[label] at (5.0, 4.1) {BMU encontrada};
% Entrada
\node[input] at (5.0, 4.55) {$\bm{x}$};
\foreach \row in {0,1,2}{
    \foreach \col in {0,1,2}{
        \pgfmathparse{int(\row==1 && \col==1)}
        \ifnum\pgfmathresult=1
            \node[bmu]      at (3.6+\col*0.9, 3.2-\row*0.9) {};
        \else
            \pgfmathparse{int(abs(\row-1)<=1 && abs(\col-1)<=1)}
            \ifnum\pgfmathresult=1
                \node[neighbor] at (3.6+\col*0.9, 3.2-\row*0.9) {};
            \else
                \node[neuron]   at (3.6+\col*0.9, 3.2-\row*0.9) {};
            \fi
        \fi
    }
}
\draw[red!60, dashed] (5.0, 4.4) -- (4.5, 2.3);
\node[font=\tiny, gray] at (5.0, 0.8) {BMU=rojo, vecinos=naranja};

% ─── Flecha paso 2 ──────────────────────────────────────────
\draw[arrow, blue!60] (6.4, 2.2) -- (7.2, 2.2);
\node[font=\tiny, align=center, blue!70] at (6.8, 2.55)
    {2. actualizar\\BMU+vecinos};

% ─── Grilla actualizada ─────────────────────────────────────
\node[label] at (9.0, 4.1) {pesos actualizados};
\foreach \row in {0,1,2}{
    \foreach \col in {0,1,2}{
        \pgfmathparse{int(\row==1 && \col==1)}
        \ifnum\pgfmathresult=1
            \node[bmu]      at (7.6+\col*0.9, 3.2-\row*0.9) {};
        \else
            \pgfmathparse{int(abs(\row-1)<=1 && abs(\col-1)<=1)}
            \ifnum\pgfmathresult=1
                \node[neighbor] at (7.6+\col*0.9, 3.2-\row*0.9) {};
            \else
                \node[neuron]   at (7.6+\col*0.9, 3.2-\row*0.9) {};
            \fi
        \fi
    }
}
\node[font=\tiny, gray, align=center] at (9.0, 0.8) {BMU y vecinos se\\acercan a $\bm{x}$};

% ─── Leyenda regla de actualización ─────────────────────────
\node[draw=blue!20, fill=blue!5, rounded corners, font=\footnotesize,
      inner sep=5pt, align=left] at (5.5, -0.5)
      {$\bm{w}_i(t{+}1) = \bm{w}_i(t) + \alpha(t)\, h_{ci}(t)\,[\bm{x}(t) - \bm{w}_i(t)]$};

\end{tikzpicture}
\caption{Proceso de entrenamiento de un SOM en una iteración. (1) Se presenta
una entrada $\bm{x}$ y se identifica la neurona ganadora o \textit{Best Matching
Unit} (BMU, en rojo), la de menor distancia a $\bm{x}$. (2) La BMU y sus
vecinas (naranja) actualizan sus pesos acercándose a $\bm{x}$ según la regla
mostrada, donde $\alpha(t)$ es la tasa de aprendizaje y $h_{ci}(t)$ es la
función de vecindad.}
\label{fig:som_proceso}
\end{figure}

% ============================================================
\begin{figure}[H]
\centering
\begin{tikzpicture}[
    neuron/.style={circle, draw=blue!50, fill=blue!8, minimum size=0.55cm},
    bmu/.style={circle, draw=red!70, fill=red!20, minimum size=0.55cm, very thick},
    neighbor/.style={circle, draw=orange!70, fill=orange!10, minimum size=0.55cm, thick},
    input/.style={circle, draw=green!60!black, fill=green!15, minimum size=0.6cm, font=\tiny},
    arrow/.style={-{Stealth}, thick},
    label/.style={font=\small\bfseries, gray}
]

% ─── Grilla 3×3 en t=0 (aleatoria) ──────────────────────────
\node[label] at (1.0, 4.1) {$t = 0$ (inicial)};
\foreach \row in {0,1,2}{
    \foreach \col in {0,1,2}{
        \node[neuron] at (\col*0.9, 3.2-\row*0.9) {};
    }
}
\node[font=\tiny, gray] at (1.0, 0.8) {pesos aleatorios};

% ─── Flecha paso 1 ──────────────────────────────────────────
\draw[arrow, blue!60] (2.4, 2.2) -- (3.2, 2.2);
\node[font=\tiny, align=center, blue!70] at (2.8, 2.55)
    {1. presentar\\entrada $\bm{x}$};

% ─── Grilla con BMU marcada ──────────────────────────────────
\node[label] at (5.0, 4.1) {BMU encontrada};
% Entrada
\node[input] at (5.0, 4.55) {$\bm{x}$};
\foreach \row in {0,1,2}{
    \foreach \col in {0,1,2}{
        \pgfmathparse{int(\row==1 && \col==1)}
        \ifnum\pgfmathresult=1
            \node[bmu]      at (3.6+\col*0.9, 3.2-\row*0.9) {};
        \else
            \pgfmathparse{int(abs(\row-1)<=1 && abs(\col-1)<=1)}
            \ifnum\pgfmathresult=1
                \node[neighbor] at (3.6+\col*0.9, 3.2-\row*0.9) {};
            \else
                \node[neuron]   at (3.6+\col*0.9, 3.2-\row*0.9) {};
            \fi
        \fi
    }
}
\draw[red!60, dashed] (5.0, 4.4) -- (4.5, 2.3);
\node[font=\tiny, gray] at (5.0, 0.8) {BMU=rojo, vecinos=naranja};

% ─── Flecha paso 2 ──────────────────────────────────────────
\draw[arrow, blue!60] (6.4, 2.2) -- (7.2, 2.2);
\node[font=\tiny, align=center, blue!70] at (6.8, 2.55)
    {2. actualizar\\BMU+vecinos};

% ─── Grilla actualizada ─────────────────────────────────────
\node[label] at (9.0, 4.1) {pesos actualizados};
\foreach \row in {0,1,2}{
    \foreach \col in {0,1,2}{
        \pgfmathparse{int(\row==1 && \col==1)}
        \ifnum\pgfmathresult=1
            \node[bmu]      at (7.6+\col*0.9, 3.2-\row*0.9) {};
        \else
            \pgfmathparse{int(abs(\row-1)<=1 && abs(\col-1)<=1)}
            \ifnum\pgfmathresult=1
                \node[neighbor] at (7.6+\col*0.9, 3.2-\row*0.9) {};
            \else
                \node[neuron]   at (7.6+\col*0.9, 3.2-\row*0.9) {};
            \fi
        \fi
    }
}
\node[font=\tiny, gray, align=center] at (9.0, 0.8) {BMU y vecinos se\\acercan a $\bm{x}$};

% ─── Leyenda regla de actualización ─────────────────────────
\node[draw=blue!20, fill=blue!5, rounded corners, font=\footnotesize,
      inner sep=5pt, align=left] at (5.5, -0.5)
      {$\bm{w}_i(t{+}1) = \bm{w}_i(t) + \alpha(t)\, h_{ci}(t)\,[\bm{x}(t) - \bm{w}_i(t)]$};

\end{tikzpicture}
\caption{Proceso de entrenamiento de un SOM en una iteración. (1) Se presenta
una entrada $\bm{x}$ y se identifica la neurona ganadora o \textit{Best Matching
Unit} (BMU, en rojo), la de menor distancia a $\bm{x}$. (2) La BMU y sus
vecinas (naranja) actualizan sus pesos acercándose a $\bm{x}$ según la regla
mostrada, donde $\alpha(t)$ es la tasa de aprendizaje y $h_{ci}(t)$ es la
función de vecindad.}
\label{fig:som_proceso}
\end{figure}
